\section{Discussion}
\label{sec:discussion}

\paragraph{When is model-response watermarking useful?}
Shibboleth is most appropriate when the production decoder should remain
unchanged and the verifier can retain the checkpoint.  It does not replace a
token-local detector: the latter is cheaper and easier to run at scale.  A
practical deployment could use a token-local test for routine checks and
reserve checkpoint replay for documents that require a model-specific
provenance claim.

\paragraph{Checkpoint binding is both the signal and the cost.}
A successful Shibboleth test says that the text agrees with keyed responses of
a particular checkpoint, not merely that its token identities follow a keyed
pattern.  This is the method's distinguishing property, but it also makes the
checkpoint identity, numerical precision, and replay implementation part of
the verification record.  Transfer across numerically different copies is not
established by our experiments.

\paragraph{What does a positive result mean?}
The detector supports a narrow statement: under the registered key and replay
procedure, the document contains more checkpoint-consistent carrier responses
than expected by chance.  It does not identify the human who requested the
text, establish that every passage came from the model, or rule out later
editing.  Those questions require account records, cryptographic signatures,
or separate forensic analysis.  Keeping this interpretation narrow is
especially important when the detector is used outside a controlled benchmark.

\paragraph{Editing reveals where synchronisation lives.}
Dispersed edits alter many later prefixes and therefore many reconstructed
probes; a contiguous span leaves all earlier carriers intact.  The measured
asymmetry suggests that content-anchored probes, rather than absolute prefix
positions, are the most direct route to editing tolerance.  The block-local
and windowed variants tested here did not improve detection after random
edits, so this remains future work rather than a demonstrated remedy.

\paragraph{Interpreting quality.}
A perplexity ratio below one does not by itself imply better text, and the task
cohorts are too small to support claims about modest accuracy changes.  The
relevant result is narrower: under the measured settings, Shibboleth's task
score changes are small relative to the degradation observed for direct
dgMARK on Dream, while its detector requires substantially more computation.

\paragraph{Reporting and deployment.}
The operational record should include the key identifier, checkpoint hash,
numerical precision, probe construction, decision threshold, and the number of
usable carriers.  Reporting the carrier count alongside the decision separates
a weak result caused by too little text from a genuinely unaligned response.
In a deployment, the verifier can also return the exact binomial probability
rather than only a binary label, leaving the downstream policy to the operator.
Our experiments validate the statistical calculation and measure several
editing patterns; they do not define such a policy or its consequences.
